
\begin{obereflection}
\textbf{Latihan 6.1}
1. Jika sebuah cipher blok berukuran 64 bit memakai mode ECB, apa yang terjadi
   bila dua blok plainteks yang identik dikirimkan?
2. Jelaskan perbedaan peran S-box dan P-box dalam mencapai \textit{confusion}
   dan \textit{diffusion}.
3. Mengapa jaringan Feistel memudahkan implementasi dekripsi?
4. Sebuah pesan berukuran 100 bit dienkripsi dengan cipher blok berukuran
   8 bit. Berapa blok yang terbentuk, dan berapa bit \textit{padding} yang
   diperlukan pada blok terakhir?
5. Pada mode CBC, sebuah bit pada blok cipherteks ke-3 rusak dalam perjalanan.
   Blok plainteks mana saja yang ikut rusak setelah dekripsi? Jelaskan.
6. Bandingkan mode ECB dan mode CTR: keduanya tidak merantai blok. Mengapa CTR
   tidak mewarisi kelemahan pola berulang yang dimiliki ECB?
7. Mengapa mode CFB dan OFB hanya memerlukan fungsi enkripsi $E_K$, baik untuk
   enkripsi maupun dekripsi?
8. Tuliskan persamaan enkripsi dan dekripsi mode CBC, lalu jelaskan peran IV.
\end{obereflection}

\begin{obereflection}
\textbf{Latihan 6.2}
1. Dengan cipher blok mainan $E_K(P)=(P\oplus K)\ll 1$ berukuran 4 bit dan
   $K=\code{1011}$, enkripsilah plainteks \code{1100 1100} dalam mode ECB, lalu
   jelaskan mengapa hasilnya memperlihatkan pengulangan.
2. Dengan cipher blok mainan yang sama dan $IV=\code{0000}$, hitung cipherteks
   dari plainteks pada soal nomor 1 memakai mode CBC.
3. Turunkan sifat reversible jaringan Feistel, yaitu bahwa
   $L_{i-1}=R_i\oplus f(R_{i-1},K_i)$ selalu mengembalikan separuh blok yang
   benar, tanpa mengasumsikan bentuk tertentu untuk $f$.
4. Jelaskan apa yang dimaksud \textit{key expansion} dan mengapa panjang kunci
   putaran tidak harus sama dengan panjang kunci eksternal.
5. Diberikan masukan 6 bit $\code{101010}$ pada S-box $S_1$ DES. Tentukan nomor
   baris dan nomor kolomnya, lalu cari keluarannya pada baris $S_1$ yang
   diberikan di dalam bab ini.
6. Sebuah cipher blok berukuran 128 bit memakai kunci 256 bit. Jelaskan
   pengaruh kedua ukuran itu terhadap diffusion dan confusion.
7. Pada mode OFB, mengapa kesalahan satu bit pada cipherteks tidak menjalar ke
   blok plainteks berikutnya?
\end{obereflection}
